% Part: proof-theory
% Chapter: natural-deduction
% Section: grafting

\documentclass[../../../include/open-logic-section]{subfiles}

\begin{document}

\olfileid{pt}{nat}{gra}

\olsection{\usetoken{P}{proof} پېوندول}

په \Log{N1c} او~\Log{N1i} کښې !!^{proof}s له فرضونو~$!B$
څخه د !!{formula}s~$!A$ !!{proof}s دي۔ فرض کړئ چې پخپله $!B$
يو !!a{proof} لري۔ د~$!B$ !!{proof} د $!A$ له هغه !!{proof}
سره څنګه يوځاے کړو چې له~$!B$ څخه دے، څو د~$!A$ داسې
يو !!a{proof} تر لاسه کړو چې فرض~$!B$ نۀ کاروي؟

يوه طريقه دا ده چې د~\Intro{\lif} په کارولو د~$!A$ له
!!{proof} څخه فرض~$!B$ وايستل شي۔ بيا د~$!B \lif !A$ نوے
!!{proof} د~$!B$ له !!{proof} سره يوځاے کولے شو چې دا تر
لاسه کړو:
\[
\AxiomC{$\Discharge{!B}{x}$}
\DeduceC{$!A$}
\DischargeRule{\Intro{\lif}}{x}
\UnaryInfC{$!B \lif !A$}
\AxiomC{}
\DeduceC{$!B$}
\RightLabel{\Elim{\lif}}
\BinaryInfC{$!A$}
\DisplayProof
\]
\textbf{سرچينه‌يي سمون OLCMP-130:} د شرط تړنې د داخلولو د
قاعدې پايله د هماغه فرض او هماغې پايلې شرط تړنه ده؛ د رسم
نااړوند توري سم شوے دے۔
دا نېغه طريقه ده، خو لږه بدرنګه: ولې $\lif$ داخل کړو چې سم
لاسي ئې بيا وايستل شي؟ د $!B \lif !A$ له لارې دا ډول
«تاو راتاو» بايد پرېښودلے شي۔ (دا چې تل پرېښودلے شي، په
حقيقت کښې د عادي کولو د قضيې مضمون دے۔)

په \Log{N1c} او~\Log{N1i} کښې په دې حالت کښې دا په اسانۍ
پرېښودلے شو: فرضونه~$!B^x$ يوازې د~$!B$ په ټول !!{proof}
بدلولے شو۔ د نورو د فرضونو پر ځاے د !!{proof}s دا «پېوندول»
ډېر ګټور دي، نو د تعريف په توګه ئې ثبتوو۔

\begin{defn}
که $\delta_1$ د پرانيستي فرض~$!B^x$ څخه د~$!A$ يو
\Log{N1c}-!!{proof} (يا \Log{N1i}-!!{proof}) وي، او $\delta_2$
د~$!B$ يو \Log{N1c}-!!{proof} (\Log{N1i}-!!{proof}) وي، نو
$\Subst{\delta_1}{\delta_2}{!B^x}$ په~$\delta_1$ کښې د~$!B^x$
د هر نۀ ايستل شوي راتګ په~$\delta_2$ د بدلولو پايله ده۔
\end{defn}

په دې شرط چې $\delta_1$ او $\delta_2$ په عين نښه د بېلابېلو
فرض !!{formula}s نۀ لري او د~$\delta_2$ هېڅ پرانيستے فرض
د~$\delta_1$ ځانګړے متغير نۀ لري، نو
$\Subst{\delta_1}{\delta_2}{!B^x}$ يو کره !!{proof} دے۔ هر
پرانيستے فرض ئې يا د $\delta_1$ يو نۀ بدل شوے پرانيستے فرض
دے، يا د~$\delta_2$ داسې پرانيستے فرض دے چې له پېوندولو وروسته
نۀ دے ايستل شوے۔ \textbf{سرچينه‌يي سمون OLCMP-133:} بدل شوي
فرضونه نۀ پاتې کېږي او نښې لرونکے داخل شوے فرض په ورپسې
قاعده کښې ايستل کېدے شي؛ د ټولو پخوانيو فرضونو د عين پاتې
کېدو دعوه سمه شوې ده۔ د !!{proof}s د پېوندولو په وخت دا شرطونه
عموماً پوره وي۔ لومړنے شرط هله پوره وي چې $\delta_1$ او
$\delta_2$ د يوه لوے !!{proof} فرعي ثبوتونه وي۔ که دويم
شرط پوره نۀ وي، د~$\delta_1$ ځانګړي متغيرونه بيا نومولے شو
چې پوره شي۔

ورته تعريف په \Log{N2c} او~\Log{N2i} هم لګېږي۔

\begin{defn}
که $\delta_1$ د~$x:!B, \Gamma_1 \Sequent !A$ يو
\Log{N2c}-!!{proof} (يا \Log{N2i}-!!{proof}) وي، او $\delta_2$
د~$\Gamma_2 \Sequent !B$ يو \Log{N2c}-!!{proof}
(\Log{N2i}-!!{proof}) وي، نو $\Subst{\delta_1}{\delta_2}{x:!B}$
په~$\delta_1$ کښې د هرې بديهيې $x:!B \Sequent !B$ په~$\delta_2$
د بدلولو او د داسې بديهياتو د لاندې هر سېکوېنټ سياق ته د
$\Gamma_2$ د زياتولو پايله ده۔
\textbf{سرچينه‌يي سمون OLCMP-131:} دويم ثبوت هم د همدغه
سېکوېنټي نظام ثبوت دے؛ د فارمولونو د بل نظام نومونه سم شول۔
\end{defn}

\begin{prop}
که $\delta_1 \Proves x:!B, \Gamma_1 \Sequent
!A$ او $\delta_2 \Proves \Gamma_2 \Sequent !B$ وي، نو
$\Subst{\delta_1}{\delta_2}{x:!B} \Proves \Gamma_1, \Gamma_2 \Sequent
!A$، په دې شرط چې د~$\delta_1$ هېڅ ځانګړے متغير په~$\Gamma_2$
کښې رانۀ شي۔
\end{prop}

\begin{proof}
د $\pheight{\delta_1}$ پر لوړوالي په استقرا۔
\textbf{سرچينه‌يي سمون OLCMP-132:} استقرا د لومړي ثبوت پر
لوړوالي ده؛ د نۀ ټاکل شوي ثبوت نوم سم شوے دے۔
\end{proof}

\end{document}
